% V2 single-column - parser/ATS-safe twin of the sidebar design.
% Single column, no multi-column boxes, no tables, no photo, real selectable text.
% Build: xelatex v2-singlecol.tex
\documentclass[10pt]{article}

\usepackage[a4paper, top=1.7cm, bottom=1.6cm, left=2.1cm, right=2.1cm]{geometry}
\usepackage{fontspec}
\usepackage{xcolor}
\usepackage{enumitem}
\usepackage[hidelinks]{hyperref}
\usepackage{microtype}
\usepackage{ragged2e}

\pagestyle{empty}
\setlength{\parindent}{0pt}
\linespread{1.06}
\setlength{\emergencystretch}{3em}
\RaggedRight

% ---- Fonts ----
\defaultfontfeatures{Path = /usr/share/fonts/rsms-inter-fonts/, Extension = .ttf, Ligatures = TeX}
\setmainfont{Inter}[UprightFont = Inter-Regular, BoldFont = Inter-SemiBold, ItalicFont = Inter-Italic]
\newfontfamily\interlight{Inter}[UprightFont = Inter-Light]
\newfontfamily\intermed{Inter}[UprightFont = Inter-Medium]
\newfontfamily\intersemi{Inter}[UprightFont = Inter-SemiBold]
\newfontfamily\displight{InterDisplay}[UprightFont = InterDisplay-Light]
\newfontfamily\dispthin{InterDisplay}[UprightFont = InterDisplay-Thin]

% ---- Palette (matches the sidebar version) ----
\definecolor{ink}{HTML}{20222A}
\definecolor{muted}{HTML}{767A85}
\definecolor{accent}{HTML}{1F5C63}
\definecolor{hair}{HTML}{D6D7DB}
\color{ink}

\setlist[itemize]{leftmargin=1.15em, itemsep=0.25ex, topsep=0.5ex, parsep=0pt,
  label=\textcolor{accent}{\fontsize{5}{5}\selectfont\textbullet}}

\newcommand{\cvsection}[1]{\vspace{2.6ex}{\intersemi\footnotesize\addfontfeature{LetterSpace=14.0}\textcolor{accent}{\MakeUppercase{#1}}}\hspace{0.85em}\raisebox{0.28ex}{\textcolor{hair}{\leaders\hrule height 0.5pt\hfill\kern0pt}}\par\vspace{1.3ex}}
\newcommand{\role}[3]{{\intersemi #1}\hfill{\small\textcolor{muted}{#3}}\par\vspace{0.12ex}{\small\textcolor{accent}{#2}}\par\vspace{0.7ex}}
\newcommand{\skillline}[2]{{\intersemi #1:}\ #2\par\vspace{0.55ex}}
\newcommand{\dotsep}{\thinspace\textcolor{hair}{\textbullet} }

\begin{document}

% ---- Header (single column, full width, no photo) ----
{\dispthin\fontsize{30}{34}\selectfont\addfontfeature{LetterSpace=2.0} Mark Hedley Jones}\par
\vspace{1.1ex}
{\intermed\large\textcolor{accent}{Robotics Perception Engineer}}\quad{\interlight\textcolor{muted}{SLAM \textbullet\ 3D mapping \textbullet\ sensor fusion}}\par
\vspace{1.3ex}
{\small\textcolor{muted}{Osaka, Japan \dotsep markhedleyjones@gmail.com \dotsep markhedleyjones.com \dotsep github.com/markhedleyjones \dotsep orcid.org/0000-0003-3387-2932}}

\cvsection{Profile}
{\interlight Over a decade building and shipping autonomous robots, from research platforms to deployed commercial fleets. PhD-qualified, spanning perception software, mechanical design and embedded hardware, with an AI-native workflow.}

\cvsection{Experience}
\role{Software Engineer {\interlight\textcolor{muted}{— mapping, simulation \& object detection}}}{SEQSENSE, Tokyo}{Nov 2018 — Present}
Perception and mapping software for SEQSENSE's autonomous robots: SQ-2, a security robot deployed as a fleet across airports, offices and public buildings in Japan, and FORRO, a delivery robot co-developed with Kawasaki Heavy Industries.
\begin{itemize}
  \item On-line lidar-based 3D mapping (point clouds / PCD) running across desktop, cloud and on-robot in real time.
  \item 3D object detection and tracking from lidar and camera; detection models trained in PyTorch and OpenVINO.
  \item Gazebo simulation environment that closely mirrors robot behaviour; robot models (URDF) generated from CAD.
\end{itemize}

\vspace{0.9ex}
\role{Post-Doctoral Researcher {\interlight\textcolor{muted}{— autonomous navigation}}}{Robotics Plus Ltd. / University of Auckland}{Dec 2017 — Oct 2018}
Navigation software for heavy-duty vehicles operating in outdoor orchard environments.
\begin{itemize}
  \item Integrated and tuned real-time SLAM, point-cloud processing and path-planning; selected and fused lidar, cameras, radar, IMU and GNSS.
  \item Personally accountable for the safety of people and infrastructure around the autonomous vehicle.
\end{itemize}

\vspace{0.9ex}
\role{Post-Doctoral Researcher {\interlight\textcolor{muted}{— robotic hardware development}}}{Robotics Plus Ltd. / University of Waikato}{Mar 2015 — Dec 2017}
Designed and built a heavy-duty outdoor vehicle, a fruit-harvesting robot and a pollination robot; led a team of post-graduate researchers.
\begin{itemize}
  \item Built ROS systems for robotic arms, precision spraying and an electric vehicle; full hardware design, integration and testing.
  \item Managed a 6–10 person research team across two universities, with budget and procurement authority.
\end{itemize}

\cvsection{Education}
\role{Doctor of Philosophy {\interlight\textcolor{muted}{— The Electrical Properties of Interfacial Double Layers}}}{University of Waikato}{2011 — 2015}
\vspace{-0.2ex}
\role{Master of Engineering {\interlight\textcolor{muted}{— first class honours, millimetre-wave power measurement}}}{University of Waikato}{2009}
\vspace{-0.2ex}
\role{Bachelor of Science {\interlight\textcolor{muted}{— electronics \& mechatronics}}}{University of Waikato}{2005 — 2008}

\cvsection{Skills}
\skillline{Robotics \& Perception}{ROS, SLAM, 3D mapping, point clouds (PCL), object detection \& tracking, Gazebo, URDF}
\skillline{Vision \& Machine Learning}{OpenCV, PyTorch, OpenVINO, camera/lidar calibration \& sensor fusion}
\skillline{Software}{C/C++ (7~years), Python (7~years), Linux (10+~years), Docker, Git}
\skillline{Design \& Fabrication}{SolidWorks (6~years), CAD/CAM, CNC, PCB \& embedded, fail-safe architecture}

\cvsection{Selected Publications}
\begin{itemize}
  \item Jones, M. H., et al. \textit{Design and Testing of a Heavy-Duty Platform for Autonomous Navigation in Kiwifruit Orchards.} Biosystems Engineering, 2019.
  \item Jones, M. H., Scott, J. \textit{Scaling of Electrode–Electrolyte Interface Model Parameters in Phosphate Buffered Saline.} IEEE Transactions on Biomedical Circuits and Systems, 2014.
  \item Williams, H., et al. (incl. M. H. Jones). \textit{Improvements to and large-scale evaluation of a robotic kiwifruit harvester.} Journal of Field Robotics, 2019.
\end{itemize}
{\small\textcolor{muted}{Full list: \href{https://markhedleyjones.com/about}{markhedleyjones.com/about}}}

\cvsection{Open Source}
\skillline{dmenu-extended}{fast file/program launcher for Linux tiling window managers; plugin system.}
\skillline{Container-Magic}{templated, reproducible containerised development environments.}
\skillline{Calibration boards}{100+ camera-calibration targets with an interactive generator; widely used in robotics.}
\skillline{Blender PCD I/O}{point-cloud import/export add-on for Blender.}

\cvsection{Awards}
\skillline{Waikato Doctoral Scholarship}{\$56,000 USD, 2010}
\skillline{Agilent Technologies Research Scholarship}{\$36,000 USD, 2009}
\skillline{ENZcon}{Best Research Paper (2011); Runner-up Best Presentation (2014)}

\cvsection{Languages \& Status}
\skillline{Nationalities}{New Zealand \& British}
\skillline{Languages}{English (native), Japanese (conversational, JLPT N3)}
\skillline{Status}{Japanese permanent residency}

\end{document}
